\chapter[Query Planner Internals]{\emj{🔍}~Query Planner Internals}

\section[¿Qué hace el Query Planner?]{\emj{🧠}~¿Qué hace el Query Planner?}

El Query Planner (o Query Optimizer) es el componente de PostgreSQL que toma una consulta SQL y decide \emph{cómo} ejecutarla físicamente.
No importa si el SQL es correcto: si el planner elige un plan malo, la consulta puede tardar segundos en lugar de milisegundos.
Entender el planner es la diferencia entre un DBA junior y uno senior.

\begin{teoriabox}{Pipeline de ejecución de una query}
\begin{enumerate}
  \item \textbf{Parser:} SQL texto → árbol sintáctico (parse tree)
  \item \textbf{Analyzer/Rewriter:} valida objetos (tablas, columnas existen), aplica reglas (vistas → subqueries)
  \item \textbf{Planner/Optimizer:} genera planes posibles, estima costos, elige el más barato
  \item \textbf{Executor:} ejecuta el plan elegido contra los datos reales
\end{enumerate}
El planner \textbf{nunca ejecuta} nada; solo genera y evalúa planes. El costo es una estimación, no un tiempo real.
\end{teoriabox}

\section[El modelo de costos]{\emj{💰}~El modelo de costos}

PostgreSQL asigna un costo a cada operación del plan basado en parámetros de configuración.
El costo no está en segundos; es una unidad abstracta donde \texttt{seq\_page\_cost = 1.0} es la referencia.

\begin{teoriabox}{Parámetros clave de costo}
\begin{itemize}
  \item \texttt{seq\_page\_cost = 1.0} — costo de leer una página secuencialmente (de disco o buffer)
  \item \texttt{random\_page\_cost = 4.0} — costo de acceso aleatorio (salto de disco). En SSD: bajar a 1.1--1.5
  \item \texttt{cpu\_tuple\_cost = 0.01} — costo de procesar una fila
  \item \texttt{cpu\_index\_tuple\_cost = 0.005} — costo de procesar una entrada de índice
  \item \texttt{cpu\_operator\_cost = 0.0025} — costo de evaluar un operador (=, <, etc.)
  \item \texttt{effective\_cache\_size} — estimación de cuánto RAM tiene el OS para caché de páginas. Afecta la decisión index scan vs seq scan
\end{itemize}
\end{teoriabox}

\begin{lstlisting}[caption={Ajustar costos para SSD}]
-- En postgresql.conf o por sesión:
SET random_page_cost = 1.1;
SET effective_cache_size = '12GB';  -- 75% de RAM total aprox

-- Verificar valores actuales:
SHOW random_page_cost;
SELECT name, setting, unit FROM pg_settings
WHERE name IN ('random_page_cost','seq_page_cost','effective_cache_size');
\end{lstlisting}

\section[Estadísticas: pg\_stats y ANALYZE]{\emj{📊}~Estadísticas: pg\_stats y ANALYZE}

El planner estima selectividad (¿cuántas filas retorna este filtro?) usando estadísticas recopiladas por \texttt{ANALYZE}.
Sin estadísticas actualizadas, el planner genera estimaciones malas y elige planes subóptimos.

\begin{lstlisting}[caption={Explorar estadísticas del planner}]
-- Estadísticas por columna:
SELECT attname, n_distinct, correlation, most_common_vals, most_common_freqs
FROM pg_stats
WHERE tablename = 'orders' AND attname = 'status';

-- n_distinct > 0: número exacto de valores únicos
-- n_distinct < 0: fracción de filas únicas (ej: -0.95 = 95% únicos)
-- correlation: 1.0 = físicamente ordenado (index scan muy eficiente)
--              0.0 = aleatorio (index scan costoso, seq scan mejor)

-- Forzar recolección de estadísticas:
ANALYZE orders;
ANALYZE VERBOSE orders;  -- muestra progreso

-- Statistics target por columna (default 100, rango 1-10000):
ALTER TABLE orders ALTER COLUMN status SET STATISTICS 500;
ANALYZE orders;
-- Mayor target = mejores estimaciones en columnas con distribución sesgada
\end{lstlisting}

\begin{tipbox}{Cuándo aumentar statistics target}
Columna con distribución muy sesgada: pocos valores dominan (ej: \texttt{status = 'completed'} es 95\% de las filas).
Con target bajo, el planner no ve los valores raros y subestima su selectividad.
Sube a 200--500 en columnas de filtros frecuentes con distribución irregular.
\end{tipbox}

\section[EXPLAIN y EXPLAIN ANALYZE]{\emj{🔬}~EXPLAIN y EXPLAIN ANALYZE}

\texttt{EXPLAIN} muestra el plan sin ejecutar. \texttt{EXPLAIN ANALYZE} ejecuta y muestra tiempos reales.
La diferencia entre filas estimadas y reales revela problemas de estadísticas.

\begin{lstlisting}[caption={Anatomía de EXPLAIN ANALYZE}]
EXPLAIN (ANALYZE, BUFFERS, FORMAT TEXT)
SELECT o.id, c.name, SUM(i.amount)
FROM orders o
JOIN customers c ON c.id = o.customer_id
JOIN items i ON i.order_id = o.id
WHERE o.created_at > NOW() - INTERVAL '30 days'
GROUP BY o.id, c.name;

-- Salida típica:
-- HashAggregate (cost=1234.56..1289.78 rows=500 width=52)
--               (actual time=45.2..47.8 rows=423 loops=1)
--   Buffers: shared hit=890 read=234
--   -> Hash Join (cost=...) (actual time=... rows=...)
--        Hash Cond: (i.order_id = o.id)
--        -> Seq Scan on items (cost=...) (actual rows=12847)
--        -> Hash (actual rows=423 loops=1)
--             -> Index Scan on orders using idx_orders_created_at
--                  Index Cond: (created_at > ...)

-- Interpretar:
-- cost=inicio..total    — estimación del planner
-- actual time=..        — tiempo real en ms
-- rows=N (estimado) vs rows=N (actual) → si difieren mucho: stats problema
-- Buffers: shared hit=X read=Y → hit=de cache, read=de disco
-- loops=N → si > 1: este nodo se repite N veces (nested loop)
\end{lstlisting}

La diferencia entre filas estimadas y reales es la señal más importante en un EXPLAIN ANALYZE.
Cuando el planner estima 5 filas pero en realidad hay 45,000, toda la cadena de decisiones (qué método de join, qué tipo de scan, si usar índice) se toma con información incorrecta.
El resultado puede ser un plan que funciona bien en staging (con datos de muestra) y explota en producción (con distribución real de datos).

\begin{lstlisting}[caption={Detectar estimaciones malas}]
-- Si "rows estimado" difiere 10x+ de "rows actual": problema de stats
-- Ejemplo de plan malo:
-- Seq Scan on orders (cost=0..50000 rows=5 width=...)
--                    (actual time=... rows=45000 ...)
-- El planner creyó que había 5 filas, había 45000 → eligió seq scan "barato"
-- pero en realidad debía haber filtrado antes.

-- Herramienta externa: https://explain.dalibo.com (visual)
-- Pega el output de EXPLAIN (FORMAT JSON) para visualización
EXPLAIN (ANALYZE, FORMAT JSON)
SELECT * FROM orders WHERE customer_id = 42;
\end{lstlisting}

\section[Métodos de acceso a datos]{\emj{📁}~Métodos de acceso a datos}

Cuando el planner necesita leer filas de una tabla, elige entre varios \emph{métodos de acceso}.
La decisión depende de cuántas filas se esperan retornar (selectividad), si existe un índice aplicable, cuántas páginas tiene la tabla, y qué tan ordenados físicamente están los datos.
Entender estos métodos es clave para interpretar un EXPLAIN y saber si el plan elegido es óptimo.

\begin{teoriabox}{Seq Scan vs Index Scan vs Index Only Scan}
\begin{itemize}
  \item \textbf{Seq Scan:} lee cada página de la tabla en orden. Costo fijo por página. Mejor cuando se leen $>$20--30\% de las filas.
  \item \textbf{Index Scan:} usa el índice para encontrar TIDs (row pointers), luego hace heap fetch por cada TID. Ideal para baja selectividad ($<$5\% de filas).
  \item \textbf{Index Only Scan:} todas las columnas requeridas están en el índice. No hay heap fetch. Más rápido que Index Scan pero requiere visibility map actualizado (VACUUM frecuente).
  \item \textbf{Bitmap Index Scan:} acumula TIDs del índice en un bitmap, luego hace heap fetches ordenados por bloque físico. Reduce I/O aleatorio. Aparece cuando el planner espera muchos hits de índice.
\end{itemize}
\end{teoriabox}

A veces es útil forzar al planner a usar un método alternativo para comparar cuál es más rápido en la práctica.
Esto \textbf{nunca debe hacerse en producción de forma permanente}: solo sirve para diagnóstico.
Si el planner elige un plan malo incluso después de ANALYZE, el problema real suele ser estadísticas desactualizadas o un \texttt{statistics target} muy bajo.

\begin{lstlisting}[caption={Forcing plan methods (debugging solamente)}]
-- Deshabilitar métodos para comparar planes:
SET enable_seqscan = OFF;
SET enable_hashjoin = OFF;
SET enable_mergejoin = OFF;

-- Ver el plan alternativo:
EXPLAIN SELECT * FROM orders WHERE status = 'pending';

-- IMPORTANTE: restaurar siempre después de diagnosticar
RESET enable_seqscan;
RESET enable_hashjoin;
\end{lstlisting}

\section[Tipos de Join]{\emj{🔗}~Tipos de Join}

Cuando una query une dos o más tablas, el planner debe elegir un \emph{algoritmo de join}.
Esta es una de las decisiones más costosas porque depende del tamaño de ambas relaciones, la disponibilidad de índices, la cantidad de \texttt{work\_mem} disponible, y si los datos ya vienen ordenados.
Un join mal elegido puede convertir una query de 50ms en una de 50 segundos.

\begin{teoriabox}{Nested Loop vs Hash Join vs Merge Join}
\begin{itemize}
  \item \textbf{Nested Loop:} para cada fila del outer, busca en el inner. Costo $O(N \times M)$ sin índice. Con índice en inner: $O(N \times \log M)$. Ideal cuando outer es pequeño y inner tiene índice.
  \item \textbf{Hash Join:} construye hash table del inner (el más pequeño), luego escanea outer. Costo $O(N + M)$. Ideal para joins grandes sin índice. Requiere \texttt{work\_mem} para el hash table.
  \item \textbf{Merge Join:} ambas relaciones ordenadas → un solo pasado. Costo $O(N \log N + M \log M)$ si hay que ordenar. Ideal cuando ambas entradas ya vienen ordenadas (por índice o SORT previo).
\end{itemize}
\end{teoriabox}

El Hash Join es el algoritmo más común para joins entre tablas grandes sin índices.
PG construye un hash table con la tabla más pequeña (inner), luego itera la tabla grande (outer) calculando el hash de cada fila para encontrar coincidencias.
El problema ocurre cuando el hash table no cabe en \texttt{work\_mem}: PG lo derrama a disco en ``batches'', multiplicando el I/O.
El síntoma aparece en EXPLAIN ANALYZE como \texttt{Batches: N} donde $N > 1$.

\begin{lstlisting}[caption={work\_mem y Hash Join}]
-- Si hash join usa disco (spill), aumentar work_mem:
SET work_mem = '256MB';

-- Verificar si hubo spill:
EXPLAIN (ANALYZE, BUFFERS)
SELECT * FROM big_table a JOIN other_big b ON a.id = b.id;
-- Buscar: "Batches: 4" → hubo spill a disco (ideal: Batches: 1)

-- CUIDADO: work_mem es por operación, por conexión.
-- 100 conexiones × 4 sort ops × 256MB = 102 GB RAM posible
-- Ajustar globalmente con cuidado; preferir por sesión/query
\end{lstlisting}

\begin{entrevistabox}
\begin{enumerate}
  \item ¿Qué información usa el Query Planner para estimar el costo de una query?
  \item ¿Cuándo el planner elige un Seq Scan sobre un Index Scan aunque exista el índice?
  \item ¿Qué significa que las filas estimadas y reales difieran 100x en EXPLAIN ANALYZE?
  \item Explica la diferencia entre Index Scan e Index Only Scan. ¿Qué se necesita para que el segundo funcione eficientemente?
  \item ¿Por qué \texttt{random\_page\_cost} es clave en servidores con SSD y cómo lo ajustarías?
  \item ¿Qué hace ANALYZE y cuándo se debe correr manualmente?
\end{enumerate}
\end{entrevistabox}
